\documentclass[12pt]{article}
\usepackage[a4paper,margin=1in]{geometry}
\usepackage[utf8]{inputenc}
\usepackage[T1]{fontenc}
\usepackage{amsmath,amssymb}
\usepackage{textcomp}
\usepackage{parskip}
\begin{document}
\begin{center}
{\Large \textbf{2019 STEP III}}
\end{center}

\subsection*{Page 1}
\begin{flushleft}
ERRATUM\par
Question 5 part (ii)\par
The sentence\par
"Use the substitution , where p is a suitably chosen constant, …"\par
𝑦\par
Should read\par
"Use the substitution , where p is a suitably chosen constant, …"\par
𝑦\par
\end{flushleft}
\subsection*{Page 2}
\begin{flushleft}
\bigskip\textbf{Section A: Pure Mathematics}\par
\bigskip\textbf{Question 1.} The coordinates of a particle at time t are x and y. For t \leq 0, they satisfy the pair of coupled\par
differential equations\par
\[x˙ = x ky\]
- -\par
\[y˙ = x y\]
-\par
\[where k is a constant. When t = 0, x = 1 and y = 0.\]
\textbf{(i)} Let k = 1. Find x and y in terms of t and sketch y as a function of t.\par
Sketch the path of the particle in the x-y plane, giving the coordinates of the point at\par
which y is greatest and the coordinates of the point at which x is least.\par
\textbf{(ii)} Instead, let k = 0. Find x and y in terms of t and sketch the path of the particle in\par
the x-y plane.\par
\bigskip\textbf{Question 2.} The definition of the derivative f of a (differentiable) function f is\par
\[\leq\]
f(x + h) f(x)\par
\[f \leq (x) = lim - . ( )\]
h 0 h ∗\par
\to\par
\textbf{(i)} The function f has derivative f and satisfies\par
\[\leq\]
\[f(x + y) = f(x)f(y)\]
\[for all x and y, and f (0) = k where k = 0 . Show that f(0) = 1.\]
\[\leq\]
\[\geq\]
\[Using ( ), show that f (x) = kf(x) and find f(x) in terms of x and k.\]
\[\leq\]
∗\par
\textbf{(ii)} The function g has derivative g and satisfies\par
\[\leq\]
g(x) + g(y)\par
\[g(x + y) =\]
1 + g(x)g(y)\par
\[for all x and y, g(x) < 1 for all x, and g (0) = k where k = 0 .\]
\[\leq\]
\[| | \geq\]
Find g (x) in terms of g(x) and k, and hence find g(x) in terms of x and k.\par
\[\leq\]
\end{flushleft}
\subsection*{Page 3}
\begin{flushleft}
\bigskip\textbf{Question 3.} The matrix A is given by\par
a b\par
\[A = .\]
\[\leq c d \geq\]
\textbf{(i)} You are given that the transformation represented by A has a line L1 of invariant\par
\[points (so that each point on L^1 is transformed to itself). Let (x, y) be a point on L^1.\]
\[Show that (a 1)(d 1) bc xy = 0.\]
- - -\par
( )\par
\[Show further that (a 1)(d 1) = bc.\]
- -\par
\[What can be said about A if L^1 does not pass through the origin?\]
\textbf{(ii)} By considering the cases b = 0 and b = 0 separately, show that if (a 1)(d 1) = bc\par
\[\geq - -\]
then the transformation represented by A has a line of invariant points. You should\par
identify the line in the different cases that arise.\par
\textbf{(iii)} You are given instead that the transformation represented by A has an invariant line L2\par
\[(so that each point on L^2 is transformed to a point on L^2) and that L^2 does not pass\]
\[through the origin. If L^2 has the form y = mx + k, show that (a 1)(d 1) = bc.\]
- -\par
\bigskip\textbf{Question 4.} The nth degree polynomial P(x) is said to be reflexive if:\par
\[(a) P(x) is of the form xn a1xn\]
-\par
\[1 + a2xn\]
-\par
2 + ( 1)na\par
n\par
\[where n \leq 1;\]
- - \cdots -\par
\[(b) a^1, a^2, . . . , a\]
n\par
are real;\par
\[(c) the n (not necessarily distinct) roots of the equation P(x) = 0 are a^1, a^2, . . . , a\]
n\par
.\par
\textbf{(i)} Find all reflexive polynomials of degree less than or equal to 3.\par
\textbf{(ii)} For a reflexive polynomial with n > 3, show that\par
2 2 2\par
\[2a2 = a\]
a\par
a\par
n\par
.\par
- - - \cdots -\par
Deduce that, if all the coefficients of a reflexive polynomial of degree n are integers\par
\[and a = 0, then n \geq 3.\]
n\par
\[\geq\]
\textbf{(iii)} Determine all reflexive polynomials with integer coefficients.\par
\end{flushleft}
\subsection*{Page 4}
\begin{flushleft}
\bigskip\textbf{Question 5.} (i) Let\par
x\par
\[f(x) = ,\]
\[x^2 + p\]
\[\int\]
\[where p is a non-zero constant. Sketch the curve y = f(x) for x \leq 0 in the case p > 0.\]
\textbf{(ii)} Let\par
\[I = dy,\]
\[\sum (b^2 y^2) c^2 y^2\]
- -\par
\[\int\]
bx\par
\[where b and c are positive constants. Use the substitution y = , where p is a\]
\[x^2 + p\]
\[suitably chosen constant, to show that \int\]
\[I = dx.\]
\[\sum b^2 + (b^2 c^2)x^2\]
-\par
Evaluate\par
\[\sqrt\{2\}\]
dy.\par
\[\sum1 (3 y^2) 2 y^2\]
- -\par
\[\int\]
1 1 x\par
\[[ Note: dx = tan + constant.]\]
\[\sum a^2 + x^2 a - a\]
Hence evaluate\par
y\par
dy.\par
\[\sum 1 (3y2 1) 2y2 1\]
\[\sqrt\{2\} - -\]
\[\int\]
\textbf{(iii)} By means of a suitable substitution, evaluate\par
dy.\par
\[\sum 1 (3y2 1) 2y2 1\]
\[\sqrt\{2\} - -\]
\[\int\]
\end{flushleft}
\subsection*{Page 5}
\begin{flushleft}
\bigskip\textbf{Question 6.} The point P in the Argand diagram is represented by the the complex number z, which\par
satisfies\par
\[zz az a z + aa r = 0.\]
∗ ∗ ∗ ∗\par
- - -\par
2 2\par
\[Here, r is a positive real number and r = a a. By writing z a as (z a)(z a) , show\]
∗ ∗\par
\[\geq | - | - -\]
that the locus of P is a circle, C, the radius and the centre of which you should give.\par
\textbf{(i)} The point Q is represented by ω, and is related to P by ω = . Let C be the locus of\par
\[\leq\]
z\par
Q. Show that C is also a circle, and give its radius and centre.\par
\[\leq\]
If C and C are the same circle, show that\par
\[\leq\]
2 2 2\par
\[( a r ) = 1\]
| | -\par
and that either a is real or a is imaginary. Give sketches to indicate the position of C\par
in these two cases.\par
\textbf{(ii)} Suppose instead that the point Q is represented by ω, where ω = . If the locus of Q\par
z\par
∗\par
is C, is it the case that either a is real or a is imaginary?\par
\bigskip\textbf{Question 7.} The Devil's Curve is given by\par
2 2 2 2 2 2\par
\[y (y b ) = x (x a ) ,\]
- -\par
where a and b are positive constants.\par
\textbf{(i)} In the case a = b , sketch the Devil's Curve.\par
\textbf{(ii)} Now consider the case a = 2 and b = \sqrt5, and x \leq 0, y \leq 0.\par
(a) Show by considering a quadratic equation in x 2 that either 0 \geq y \geq 1 or y \leq 2 .\par
(b) Describe the curve very close to and very far from the origin.\par
(c) Find the points at which the tangent to the curve is parallel to the x-axis and\par
the point at which the tangent to the curve is parallel to the y-axis.\par
Sketch the Devil's Curve in this case.\par
\textbf{(iii)} Sketch the Devil's Curve in the case a = 2 and b = \sqrt5 again, but with < x <\par
-\infty \infty\par
\[and < y < .\]
-\infty \infty\par
\end{flushleft}
\subsection*{Page 6}
\begin{flushleft}
\bigskip\textbf{Question 8.} A pyramid has a horizontal rectangular base ABCD and its vertex V is vertically above the\par
centre of the base. The acute angle between the face AVB and the base is α, the acute angle\par
between the face BVC and the base is β and the obtuse angle between the faces AVB and\par
\[BVC is \pi θ.\]
-\par
\textbf{(i)} The edges AB and BC are parallel to the unit vectors i and j, respectively, and the\par
unit vector k is vertical. Find a unit vector that is perpendicular to the face AVB.\par
Show that\par
\[cos θ = cos α cos β .\]
\textbf{(ii)} The edge BV makes an angle φ with the base. Show that\par
2 2 2\par
\[cot φ = cot α + cot β.\]
Show also that\par
2 2 2 2\par
cos α + cos β 2 cos θ 2 cos θ 2 cos θ\par
\[cos 2 φ = - \leq -\]
\[1 cos2 θ 1 cos2 θ\]
- -\par
\[and deduce that φ < θ .\]
\end{flushleft}
\subsection*{Page 7}
\begin{flushleft}
\bigskip\textbf{Section B: Mechanics}\par
\bigskip\textbf{Question 9.} In this question, i and j are perpendicular unit vectors and j is vertically upwards.\par
A smooth hemisphere of mass M and radius a rests on a smooth horizontal table with its\par
plane face in contact with the table. The point A is at the top of the hemisphere and the\par
point O is at the centre of its plane face.\par
Initially, a particle P of mass m rests at A. It is then given a small displacement in the\par
positive i direction. At a later time t, when the particle is still in contact with the hemisphere,\par
\[the hemisphere has been displaced by s i and ∠AOP = θ.\]
-\par
\textbf{(i)} Let r be the position vector of the particle at time t with respect to the initial position\par
of O. Write down an expression for r in terms of a, θ and s and show that\par
\[r˙ = (aθ˙ cos θ s˙) i aθ˙ sin θ j.\]
- -\par
Show also that\par
\[s˙ = (1 k) a θ˙ cos θ,\]
-\par
M\par
\[where k = , and deduce that\]
m + M\par
\[r˙ = aθ˙(k cos θ i sin θ j).\]
-\par
\textbf{(ii)} Show that\par
aθ˙2\par
k cos\par
θ + sin\par
\[θ = 2g(1 cos θ).\]
-\par
( )\par
\textbf{(iii)} At time T , when θ = α, the particle leaves the hemisphere. By considering the\par
component of ¨r parallel to the vector sin θ i + k cos θ j, or otherwise, show that at\par
time T\par
aθ˙2\par
\[= g cos α.\]
\[Find a cubic equation for cos α and deduce that cos α > .\]
\end{flushleft}
\subsection*{Page 8}
\begin{flushleft}
\bigskip\textbf{Question 10.} Two identical smooth spheres P and Q can move on a smooth horizontal table. Initially,\par
P moves with speed u and Q is at rest. Then P collides with Q. The direction of travel of P\par
before the collision makes an acute angle α with the line joining the centres of P and Q at\par
the moment of the collision. The coefficient of restitution between P and Q is e where e < 1.\par
As a result of the collision, P has speed v and Q has speed w, and P is deflected through an\par
angle θ.\par
\textbf{(i)} Show that\par
\[u sin α = v sin(α + θ)\]
and find an expression for w in terms of v, θ and α.\par
\textbf{(ii)} Show further that\par
\[sin θ = cos(θ + α) sin α + e sin(θ + α) cos α\]
and find an expression for tan θ in terms of tan α and e.\par
Find, in terms of e, the maximum value of tan θ as α varies.\par
\end{flushleft}
\subsection*{Page 9}
\begin{flushleft}
\bigskip\textbf{Section C: Probability and Statistics}\par
\bigskip\textbf{Question 11.} The number of customers arriving at a builders' merchants each day follows a Poisson distri-\par
bution with mean λ. Each customer is offered some free sand. The probability of any given\par
customer taking the free sand is p.\par
\textbf{(i)} Show that the number of customers each day who take sand follows a Poisson distri-\par
bution with mean pλ.\par
\textbf{(ii)} The merchant has a mass S of sand at the beginning of the day. Each customer who\par
\[takes the free sand gets a proportion k of the remaining sand, where 0 \geq k < 1. Show\]
that by the end of the day the expected mass of sand taken is\par
1 e kpλ S.\par
-\par
-\par
( )\par
\textbf{(iii)} At the beginning of the day, the merchant's bag of sand contains a large number of\par
grains, exactly one of which is made from solid gold. At the end of the day, the\par
merchant's assistant takes a proportion k of the remaining sand. Find the probability\par
\[that the assistant takes the golden grain. Comment on the case k = 0 and on the\]
limit k 1.\par
\to\par
\[In the case pλ > 1 find the value of k which maximises the probability that the assistant\]
takes the golden grain.\par
\bigskip\textbf{Question 12.} The set S is the set of all integers from 1 to n. The set T is the set of all distinct subsets\par
of S, including the empty set ∅ and S itself. Show that T contains exactly 2n sets.\par
\[The sets A^1, A^2, . . ., A\]
m\par
, which are not necessarily distinct, are chosen randomly and\par
independently from T , and for each k (1 \geq k \geq m), the set A is equally likely to be any of\par
k\par
the sets in T .\par
\textbf{(i)} Write down the value of P(1 A1) .\par
\in\par
\textbf{(ii)} By considering each integer separately, show that P(A1 A2 = ∅) = ( 3\par
) n .\par
\cap\par
\[Find P(A^1 A^2 A^3 = ∅) and P(A^1 A^2 A\]
m\par
\[= ∅) .\]
\cap \cap \cap \cap \cdots \cap\par
\textbf{(iii)} Find P(A1 A2) , P(A1 A2 A3) and P(A1 A2 A\par
m\par
) .\par
⊆ ⊆ ⊆ ⊆ ⊆ \cdots ⊆\par
\end{flushleft}
\end{document}
